Un pendule simple est constitué d’une masse de $\mhsqHuitMasse$ et d’un fil d’une longueur de $\mhsqHuitLongueur$. À $t = 0\text{ s}$, le pendule est à un angle de $\mhsqHuitAngle$ et sa vitesse tangentielle est de $\mhsqHuitVit$, s’éloignant du centre. Trouvez :

\begin{enumerate}[label=\alph*)]
    \item L’amplitude d’oscillation du pendule, en mètres.
    \item L’énergie mécanique du pendule.
    \item La vitesse maximale atteinte par ce pendule.
\end{enumerate}